\begin{figure}[ht]
\centering
\begin{tikzpicture}[x=0.62cm, y=0.55cm, font=\small]

% Plot log10(|x|^n / n!) vs n for x = 2, 5, 10.
% y in log10; range about [-12, 4]. Map y_cm = (log10val + 12)/16 * height.
% We'll plot raw log10 values, axis spanning -12..4 (16 decades).
% Use a coordinate: x-axis n=0..16, y-axis log10 from -12 (bottom) to 4 (top).

% Axes
\draw[->, thick] (-0.3, -12) -- (16.5, -12) node[right] {$n$};
\draw[->, thick] (-0.3, -12.3) -- (-0.3, 4.8) node[above] {$\log_{10}(|x|^{n}/n!)$};

% y ticks every 4 decades
\foreach \y in {-12,-8,-4,0,4} {
  \draw (-0.45, \y) -- (-0.15, \y) node[left, xshift=-0.1cm] {$10^{\y}$};
}
% x ticks
\foreach \n in {0,4,8,12,16} {
  \draw (\n, -12.15) -- (\n, -11.85) node[below, yshift=-0.1cm] {\n};
}

% x = 2: log10(2^n/n!), exact
\draw[thick, engblack, mark=triangle*, mark size=1pt] plot coordinates {
  (0,0) (2,0.301) (4,-0.176) (6,-1.051) (8,-2.197) (10,-3.549)
  (12,-5.068) (14,-6.726) (16,-8.504) };
\node[engblack, anchor=west] at (16.1, -8.504) {$x=2$};

% x = 5: log10(5^n/n!), exact, peak near n=4..5
\draw[thick, engblue, mark=*, mark size=1pt] plot coordinates {
  (0,0) (2,1.097) (4,1.416) (6,1.336) (8,0.986) (10,0.430)
  (12,-0.293) (14,-1.155) (16,-2.137) };
\node[engblue, anchor=west] at (16.1, -2.137) {$x=5$};

% x = 10: log10(10^n/n!), exact, peak near n=9..10
\draw[thick, engred, mark=square*, mark size=1pt] plot coordinates {
  (0,0) (2,1.699) (4,2.620) (6,3.143) (8,3.394) (10,3.440)
  (12,3.320) (14,3.060) (16,2.679) };
\node[engred, anchor=west] at (16.1, 2.679) {$x=10$};

\end{tikzpicture}
\caption[Term magnitudes of the exponential series peak then collapse]{Term
magnitudes $\abs{x}^{n}/n!$ of the exponential series for $x = 2, 5,
10$, on a log scale. Each curve rises while $n < \abs{x}$, peaks near
$n \approx \abs{x}$, and then falls faster than any geometric sequence
because the consecutive-term ratio $\abs{x}/(n+1)$ keeps shrinking. The
post-peak collapse is the factorial winning the growth hierarchy of
section 4.1, and it is why a fixed-accuracy truncation of the
exponential, sine, or cosine series needs only a handful of terms once
$\abs{x}$ is modest.}
\label{fig:vol02:ch04:factorial-decay}
\end{figure}
